% Part: first-order-logic
% Chapter: sequent-calculus
% Section: derivations

\documentclass[../../../include/open-logic-section]{subfiles}

\begin{document}

\iftag{FOL}
      {\olfileid[hi]{fol}{seq}{der}}
      {\olfileid[hi]{pl}{seq}{der}}

\olsection{\usetoken{p}{derivation}}

\begin{explain}
हम प्रारंभिक अनुक्रम का रूप और अनुमान-नियम दे चुके हैं। अनुक्रम-कलन की
!!^{derivation}s इन्हीं से आगमनात्मक रूप से बनती हैं: हर !!{derivation} या तो
अकेला प्रारंभिक अनुक्रम है, या उसका रूप है: एक अथवा दो !!{derivation}s और
उनके बाद आने वाला एक अनुमान।
\end{explain}

\begin{defn}[$\Log{LK}$ !!{derivation}]
\emph{$\Log{LK}$-!!{derivation}} अनुक्रम~$S$ के अनुक्रमों का ऐसा परिमित वृक्ष
है जो निम्नलिखित शर्तें पूरी करे:
\begin{enumerate}
\item वृक्ष के शीर्षस्थ अनुक्रम प्रारंभिक अनुक्रम हैं।
\item वृक्ष का सबसे निचला अनुक्रम~$S$ है।
\item वृक्ष में $S$ को छोड़कर हर अनुक्रम किसी अनुमान-नियम के सही प्रयोग का
  पूर्वाधार है, और उस प्रयोग का निष्कर्ष वृक्ष में उस अनुक्रम के ठीक नीचे है।
\end{enumerate}
तब $S$ को !!{derivation} का \emph{अंतिम अनुक्रम} कहते हैं और कहते हैं कि $S$,
\emph{$\Log{LK}$ में !!{derivable}} है (या $\Log{LK}$-!!{derivable} है)।
\end{defn}

\begin{ex}
हर प्रारंभिक अनुक्रम, जैसे $!C \Sequent !C$, !!a{derivation} है। इस पर, मान
लें, $\LeftR{\Weakening}$ नियम लगाकर नई !!{derivation} मिलती है:
\begin{prooftree}
\Axiom$ \Gamma \fCenter \Delta $
\RightLabel{\LeftR{\Weakening}}
\UnaryInf$ !A, \Gamma \fCenter \Delta$
\end{prooftree}
पर यह नियम सामान्य है: इसमें $!A$ के स्थान पर कोई भी !!{sentence}, जैसे~$!D$,
रख सकते हैं। यदि पूर्वाधार हमारे प्रारंभिक अनुक्रम $!C \Sequent !C$ से मेल
खाए, तो $\Gamma$ और $\Delta$ दोनों केवल~$!C$ हैं और निष्कर्ष
$!D, !C \Sequent !C$ होगा। अतः यह !!a{derivation} है:
\begin{prooftree}
\Axiom$ !C \fCenter !C $
\RightLabel{\LeftR{\Weakening}}
\UnaryInf$ !D, !C \fCenter !C$
\end{prooftree}
अब $\LeftR{\Exchange}$ जैसा दूसरा नियम लगा सकते हैं, जो बाईं ओर के दो
!!{sentence}s की अदला-बदली करता है। अतः यह भी सही !!{derivation} है:
\begin{prooftree}
\Axiom$ !C \fCenter !C $
\RightLabel{\LeftR{\Weakening}}
\UnaryInf$ !D, !C \fCenter !C$
\RightLabel{\LeftR{\Exchange}}
\UnaryInf$ !C, !D \fCenter !C$
\end{prooftree}
इस प्रयोग में नियम का दिया हुआ रूप था:
\begin{prooftree}
\Axiom$ \Gamma, !A, !B, \Pi \fCenter \Delta $
\RightLabel{\LeftR{\Exchange}}
\UnaryInf$ \Gamma, !B, !A, \Pi \fCenter \Delta,$
\end{prooftree}
यहाँ $\Gamma$ और $\Pi$ दोनों रिक्त थे, $\Delta$ का मान $!C$ था, और $!A$ तथा
$!B$ की भूमिकाएँ क्रमशः $!D$ और~$!C$ ने निभाईं। लगभग इसी तरह हम देखते हैं कि
\begin{prooftree}
\Axiom$ !D \fCenter !D $
\RightLabel{\LeftR{\Weakening}}
\UnaryInf$ !C, !D \fCenter !D$
\end{prooftree}
!!a{derivation} है। अब ये दोनों !!{derivation}s $\RightR{\land}$ से जोड़ी जा
सकती हैं। वह नियम था:
\begin{prooftree}
\Axiom$\Gamma \fCenter \Delta, !A$
\Axiom$ \Gamma \fCenter \Delta, !B$
\RightLabel{\RightR{\land}}
\BinaryInf$ \Gamma \fCenter \Delta, !A \land !B $
\end{prooftree}
हमारे मामले में !!{derivation}s पूर्वाधारों पर समाप्त होती हैं, और उनके अंतिम
अनुक्रमों का मेल उन्हीं पूर्वाधारों से होना चाहिए। अर्थात् $\Gamma$ का मान
$!C, !D$ है, $\Delta$ रिक्त है, $!A$ का मान $!C$ और $!B$ का मान $!D$ है।
अतः अनुमान सही होने पर निष्कर्ष $!C, !D \Sequent !C
\land !D$ है।
\begin{prooftree}
\Axiom$ !C \fCenter !C $
\RightLabel{\LeftR{\Weakening}}
\UnaryInf$ !D, !C \fCenter !C$
\RightLabel{\LeftR{\Exchange}}
\UnaryInf$ !C, !D \fCenter !C$
\Axiom$ !D \fCenter !D $
\RightLabel{\LeftR{\Weakening}}
\UnaryInf$ !C, !D \fCenter !D$
\RightLabel{\RightR{\land}}
\BinaryInf$ !C, !D \fCenter !C \land !D $
\end{prooftree}
पूर्वाधारों का क्रम उलट भी सकते हैं; तब $!A$ का मान $!D$ और $!B$ का मान~$!C$
होगा।
\begin{prooftree}
\Axiom$ !D \fCenter !D $
\RightLabel{\LeftR{\Weakening}}
\UnaryInf$ !C, !D \fCenter !D$
\Axiom$ !C \fCenter !C $
\RightLabel{\LeftR{\Weakening}}
\UnaryInf$ !D, !C \fCenter !C$
\RightLabel{\LeftR{\Exchange}}
\UnaryInf$ !C, !D \fCenter !C$
\RightLabel{\RightR{\land}}
\BinaryInf$ !C, !D \fCenter !D \land !C $
\end{prooftree}
\end{ex}

\end{document}
